\documentclass[12pt]{article}
\usepackage[utf8]{inputenc}
\usepackage[spanish]{babel}
\usepackage{amsmath, amssymb}
\usepackage{booktabs}
\usepackage{geometry}
\usepackage{xcolor}
\usepackage{mdframed}
\usepackage{enumitem}
\usepackage{array}
\usepackage{listings}
\geometry{margin=2.5cm}

% Cajas de color para conceptos clave
\newmdenv[
  backgroundcolor=blue!8,
  linecolor=blue!40,
  linewidth=1.5pt,
  innertopmargin=8pt,
  innerbottommargin=8pt,
  innerleftmargin=10pt,
  innerrightmargin=10pt
]{cajazul}

\newmdenv[
  backgroundcolor=orange!10,
  linecolor=orange!50,
  linewidth=1.5pt,
  innertopmargin=8pt,
  innerbottommargin=8pt,
  innerleftmargin=10pt,
  innerrightmargin=10pt
]{cajanota}

\newmdenv[
  backgroundcolor=green!8,
  linecolor=green!50,
  linewidth=1.5pt,
  innertopmargin=8pt,
  innerbottommargin=8pt,
  innerleftmargin=10pt,
  innerrightmargin=10pt
]{cajaverde}

\newmdenv[
  backgroundcolor=red!8,
  linecolor=red!40,
  linewidth=1.5pt,
  innertopmargin=8pt,
  innerbottommargin=8pt,
  innerleftmargin=10pt,
  innerrightmargin=10pt
]{cajaroja}

% Estilo para código AMPL
\lstdefinestyle{ampl}{
  basicstyle=\ttfamily\small,
  keywordstyle=\bfseries\color{blue!70!black},
  commentstyle=\color{green!50!black},
  numbers=left,
  numberstyle=\tiny\color{gray},
  frame=single,
  breaklines=true,
  morekeywords={set, param, var, minimize, maximize, s.t., sum, integer},
}

\title{\textbf{Modelado de Programaci\'on Lineal Entera-Mixta} \\[0.5em]
\large Problema 2.10 --- Problema de la Dieta --- Gu\'ia Pedag\'ogica Completa}
\author{}
\date{}

\begin{document}
\maketitle
\tableofcontents
\newpage

%% =====================================================================
\section{Antes de Modelar: \textquestiondown Qu\'e tipo de problema es este?}
%% =====================================================================

El problema de la dieta es uno de los problemas cl\'asicos m\'as antiguos de la programaci\'on lineal (fue planteado por George Stigler en 1945). La idea es simple: quieres alimentarte cumpliendo ciertos requerimientos nutricionales, gastando lo m\'inimo posible.

\begin{cajazul}
\textbf{Diferencia clave con el problema de Panam\'a/Viajeros.} En esos problemas quer\'iamos \textbf{maximizar} una ganancia. Aqu\'i queremos \textbf{minimizar} un costo. Esto cambia la direcci\'on de la optimizaci\'on, pero la estructura del modelado es la misma.

Adem\'as, las restricciones cambian de sentido:
\begin{itemize}
    \item En Panam\'a: ``no puedes gastar \textbf{m\'as de} X'' $\Rightarrow$ restricciones $\leq$
    \item En la dieta: ``debes consumir \textbf{al menos} X'' $\Rightarrow$ restricciones $\geq$
\end{itemize}
\end{cajazul}

Otra diferencia importante: este problema \textbf{no tiene \'indice de personas}. Solo hay una persona (t\'u) decidiendo qu\'e comer. Toda la estructura gira en torno a dos dimensiones: \textbf{alimentos} y \textbf{nutrientes}. Pero dentro de los alimentos hay una sutileza que lo convierte en un \textbf{problema entero-mixto (PLEM)}, como veremos.

%% =====================================================================
\section{El Problema: Descripci\'on Completa}
%% =====================================================================

\textbf{Escenario:} Debido a la recesi\'on econ\'omica pos Covid-19, usted ha decidido seguir una dieta estricta y econ\'omica que logre aportar el requerimiento m\'inimo diario nutricional para sobrevivir, y as\'i ahorrar para empezar un emprendimiento.

Asumiendo que usted es una mujer con actividad f\'isica moderada, mayor de 18 a\~nos y con un peso de 60 kg, su necesidad nutricional promedio diaria es:

\begin{table}[h]
\centering
\begin{tabular}{lc}
\toprule
\textbf{Macronutriente o micronutriente} & \textbf{Requerimiento m\'inimo} \\
\midrule
Prote\'ina (g/d\'ia)           & 65    \\
L\'ipidos o grasas (g/d\'ia)   & 66    \\
Hidratos de carbono (g/d\'ia)  & 362   \\
Vitamina A ($\mu$g/d\'ia)      & 700   \\
Vitamina C (mg/d\'ia)          & 75    \\
Folato (mg/d\'ia)              & 400   \\
Calcio (mg/d\'ia)              & 1\,000 \\
Hierro (mg/d\'ia)              & 18    \\
\bottomrule
\end{tabular}
\caption{Requerimiento m\'inimo nutricional diario (Tabla 2.9 del libro).}
\end{table}

Los alimentos disponibles, su aporte nutricional por porci\'on y su costo son:

\begin{table}[h]
\centering
\small
\begin{tabular}{l *{9}{r}}
\toprule
& \textbf{LE} & \textbf{H} & \textbf{P} & \textbf{A}
& \textbf{F}  & \textbf{Z} & \textbf{M} & \textbf{N} & \textbf{Mant} \\
\midrule
Prote\'ina (g)          & 3,3  & 13,5  & 20,6 & 6,4  & 12,2 & 0,9   & 0,3  & 0,7  & 0    \\
L\'ipidos (g)           & 3,2  & 10    & 1,6  & 0,8  & 0,3  & 0,5   & 0,3  & 0,3  & 82,9 \\
Hidr.\ carbono (g)      & 4,8  & 4     & 57,3 & 79,7 & 74,6 & 8,1   & 14,5 & 8,7  & 0    \\
Vit.\ A ($\mu$g)        & 91   & 226   & 85   & 0    & 0    & 1\,346 & 4   & 40   & 5,6  \\
Vit.\ C (mg)            & 4,5  & 0     & 26   & 0    & 0    & 6     & 10   & 50   & 0    \\
Folato (mg)             & 14   & 51,2  & 88   & 12   & 4    & 10    & 5    & 37   & 3    \\
Calcio (mg)             & 123  & 100   & 260  & 14   & 24   & 41    & 6    & 36   & 0    \\
Hierro (mg)             & 0,1  & 2,4   & 2    & 0,8  & 3,6  & 0,7   & 0,4  & 0,3  & 0    \\
\midrule
\textbf{Costo (COP)}   & 660  & 1\,180 & 2\,100 & 950 & 560 & 440 & 1\,440 & 400 & 800 \\
\bottomrule
\end{tabular}
\caption{Aporte nutricional y costo por porci\'on (Tabla 2.10 del libro).}
\end{table}

Donde las porciones son: LE = leche entera ($\frac{1}{2}$ taza), H = huevo (2 und), P = porotos (2 tazas), A = arroz (2 tazas), F = fideos (2 tazas), Z = zanahoria (1 und), M = manzana (1 und), N = naranja (1 und), Mant = mantequilla (4 cucharadas).

\begin{cajanota}
\textbf{La pregunta que queremos responder:} \textquestiondown Cu\'antas porciones de cada alimento debo consumir al d\'ia para cubrir todos mis requerimientos nutricionales al menor costo posible?
\end{cajanota}

%% =====================================================================
\section{Paso 1 --- Identificar los Conjuntos}
%% =====================================================================

La primera pregunta de todo modelado es: \textbf{\textquestiondown qu\'e dimensiones tiene este problema?} Es decir, \textquestiondown sobre qu\'e cosas necesito indexar?

\subsection*{La separaci\'on clave: alimentos continuos vs.\ discretos}

Observa las unidades de medida de cada alimento:

\begin{table}[h]
\centering
\begin{tabular}{lll}
\toprule
\textbf{Alimento} & \textbf{Porci\'on} & \textbf{\textquestiondown Divisible?} \\
\midrule
Leche entera   & $\frac{1}{2}$ taza   & \textbf{S\'i} --- puedes tomar 1,3 medias-tazas \\
Porotos        & 2 tazas              & \textbf{S\'i} --- puedes cocinar 0,5 porciones \\
Arroz          & 2 tazas              & \textbf{S\'i} --- puedes servir 3,8 porciones \\
Fideos         & 2 tazas              & \textbf{S\'i} --- puedes cocinar fracciones \\
Mantequilla    & 4 cucharadas         & \textbf{S\'i} --- puedes usar 0,6 porciones \\
\midrule
Huevo          & 2 unidades           & \textbf{No} --- no puedes comprar 1,7 porciones de huevo \\
Zanahoria      & 1 unidad             & \textbf{No} --- compras zanahorias enteras \\
Manzana        & 1 unidad             & \textbf{No} --- compras manzanas enteras \\
Naranja        & 1 unidad             & \textbf{No} --- compras naranjas enteras \\
\bottomrule
\end{tabular}
\end{table}

\begin{cajazul}
\textbf{Esto obliga a separar los alimentos en dos conjuntos:}

\begin{itemize}
    \item $I = \{1, 3, 4, 5, 9\}$ --- alimentos \textbf{continuos}: leche, porotos, arroz, fideos, mantequilla. Sus cantidades pueden ser fraccionarias (valores reales).

    \item $A = \{2, 6, 7, 8\}$ --- alimentos \textbf{discretos}: huevo, zanahoria, manzana, naranja. Sus cantidades deben ser enteras (no puedes comprar media naranja).
\end{itemize}

N\'otese que los n\'umeros del conjunto \textbf{no son consecutivos}: $I = \{1,3,4,5,9\}$, no $\{1,2,3,4,5\}$. Esto es porque la numeraci\'on original de los alimentos se preserva (1=leche, 2=huevo, 3=porotos, etc.) y cada alimento va al conjunto que le corresponde seg\'un su naturaleza.
\end{cajazul}

\begin{cajaroja}
\textbf{Error com\'un: tratar todos los alimentos como continuos.} Si hicieras eso, el solver podr\'ia decirte ``come 0,3 naranjas al d\'ia''. \textquestiondown Qu\'e haces con 0,3 naranjas? No tiene sentido pr\'actico. La separaci\'on en $I$ y $A$ evita este problema.

Este detalle es lo que convierte al problema de \textbf{PL puro} en un \textbf{PLEM} (Programa Lineal Entero-Mixto): tiene variables continuas \textit{y} enteras.
\end{cajaroja}

Adem\'as, necesitamos un conjunto para los nutrientes:

\begin{cajazul}
\textbf{Conjunto de nutrientes} (lo que debes satisfacer):
\[
    J = \{1, 2, \ldots, 8\}
\]
para prote\'ina, l\'ipidos, hidratos de carbono, vitamina A, vitamina C, folato, calcio y hierro, respectivamente.
\end{cajazul}

\begin{cajanota}
\textbf{Resumen de conjuntos.} En total hay \textbf{tres} conjuntos: $I$ (5 alimentos continuos), $A$ (4 alimentos discretos), y $J$ (8 nutrientes). Los alimentos est\'an en $I \cup A$ (la uni\'on de ambos), que tiene 9 elementos.
\end{cajanota}

%% =====================================================================
\section{Paso 2 --- Identificar los Par\'ametros}
%% =====================================================================

Antes de definir variables, conviene nombrar los \textbf{datos} (valores conocidos) del problema. Esto es clave para distinguir entre lo que \textit{conoces} (par\'ametros) y lo que \textit{decides} (variables).

\begin{cajazul}
\textbf{Par\'ametros del problema:}
\begin{itemize}
    \item $c_i$ : costo del alimento $i \in I \cup A$.
    \item $a_{ij}$ : aporte nutricional del alimento $i \in I \cup A$ en el nutriente $j \in J$.
    \item $d_j$ : requerimiento m\'inimo del nutriente o vitamina $j \in J$ en el cuerpo humano.
\end{itemize}
\end{cajazul}

Por ejemplo:
\begin{itemize}
    \item $a_{2,1} = 13{,}5$ significa que 1 porci\'on del alimento 2 (huevo) aporta 13,5 g del nutriente 1 (prote\'ina).
    \item $d_7 = 1\,000$ significa que necesitas al menos 1\,000 mg del nutriente 7 (calcio) al d\'ia.
    \item $c_1 = 660$ significa que una porci\'on del alimento 1 (leche) cuesta 660 COP.
\end{itemize}

\begin{cajaverde}
\textbf{Regla pr\'actica.} En la Tabla 2.10, cada \textbf{columna} es un alimento $i \in I \cup A$, cada \textbf{fila} (excepto la \'ultima) es un nutriente $j \in J$, y cada \textbf{celda} es el par\'ametro $a_{ij}$. La \'ultima fila son los costos $c_i$. La Tabla 2.9 contiene los $d_j$.
\end{cajaverde}

%% =====================================================================
\section{Paso 3 --- Definir las Variables de Decisi\'on}
%% =====================================================================

Aqu\'i viene el punto central que distingue a este problema. Como los alimentos se dividieron en dos conjuntos, necesitamos \textbf{dos tipos de variables}:

\begin{cajazul}
\textbf{Variables de decisi\'on:}
\begin{itemize}
    \item $x_i$ : cantidad de porciones del alimento $i \in I$ a comprar diariamente. \textbf{Continua.}
    \item $y_i$ : cantidad de porciones del alimento $i \in A$ a comprar diariamente. \textbf{Entera.}
\end{itemize}
\end{cajazul}

Es decir, tenemos:

\begin{table}[h]
\centering
\begin{tabular}{llll}
\toprule
\textbf{Variable} & \textbf{Alimento} & \textbf{Tipo} & \textbf{Dominio} \\
\midrule
$x_1$ & Leche entera   & Continua & $\mathbb{R}^+ \cup \{0\}$ \\
$x_3$ & Porotos        & Continua & $\mathbb{R}^+ \cup \{0\}$ \\
$x_4$ & Arroz          & Continua & $\mathbb{R}^+ \cup \{0\}$ \\
$x_5$ & Fideos         & Continua & $\mathbb{R}^+ \cup \{0\}$ \\
$x_9$ & Mantequilla    & Continua & $\mathbb{R}^+ \cup \{0\}$ \\
\midrule
$y_2$ & Huevo          & Entera   & $\mathbb{Z}^+ \cup \{0\}$ \\
$y_6$ & Zanahoria      & Entera   & $\mathbb{Z}^+ \cup \{0\}$ \\
$y_7$ & Manzana        & Entera   & $\mathbb{Z}^+ \cup \{0\}$ \\
$y_8$ & Naranja        & Entera   & $\mathbb{Z}^+ \cup \{0\}$ \\
\bottomrule
\end{tabular}
\caption{Variables de decisi\'on del problema de la dieta.}
\end{table}

\begin{cajaroja}
\textbf{\textquestiondown Por qu\'e no una sola variable $x_i$ para todos?} Porque si usaras una sola variable continua para todos, el solver podr\'ia asignar $x_2 = 0{,}3$ (comprar 0,3 porciones de huevo), lo cual no tiene sentido f\'isico. Y si las hicieras todas enteras, perder\'ias la flexibilidad de comprar 4,071 medias-tazas de leche (que s\'i es posible).

La separaci\'on en $x$ (continuas) e $y$ (enteras) captura correctamente la realidad del problema. Esto convierte el problema en un \textbf{PLEM}: Programa Lineal Entero-Mixto.
\end{cajaroja}

\begin{cajanota}
\textbf{An\'alisis dimensional.} Las variables $x$ e $y$ se miden en \textbf{porciones por d\'ia}. El costo $c_i$ se mide en COP por porci\'on. Al multiplicar: $c_i \cdot x_i$ da COP/d\'ia, que es la unidad correcta para el costo diario. Esto se puede verificar as\'i:

\[
    \left[\frac{\text{COP}}{\text{d}}\right]
    = \left[\frac{\text{COP}}{\text{continuo}}\right] \cdot \left[\frac{\text{continuo}}{\text{d}}\right]
    + \left[\frac{\text{COP}}{\text{discreto}}\right] \cdot \left[\frac{\text{discreto}}{\text{d}}\right]
\]
\end{cajanota}

%% =====================================================================
\section{Paso 4 --- Funci\'on Objetivo: \textquestiondown Qu\'e minimizamos?}
%% =====================================================================

Queremos \textbf{minimizar el costo total diario} de la dieta. Como hay dos tipos de variables, la funci\'on objetivo tiene \textbf{dos sumatorias}: una sobre los alimentos continuos y otra sobre los discretos.

Expl\'icitamente:
\begin{align*}
    \text{Min}\; z &= \underbrace{660\,x_1 + 2\,100\,x_3 + 950\,x_4 + 560\,x_5 + 800\,x_9}_{\text{alimentos continuos } (i \in I)} \\[4pt]
    &+ \underbrace{1\,180\,y_2 + 440\,y_6 + 1\,440\,y_7 + 400\,y_8}_{\text{alimentos discretos } (i \in A)}
\end{align*}

En notaci\'on algebraica compacta:

\[
    \text{Min}\; z = \sum_{i \in I} c_i \cdot x_i + \sum_{i \in A} c_i \cdot y_i
\]

\begin{cajazul}
\textbf{\textquestiondown Por qu\'e dos sumatorias y no una?} Porque los alimentos continuos usan la variable $x_i$ y los discretos usan $y_i$. Son variables diferentes. Si todos fueran del mismo tipo, bastar\'ia una sola sumatoria $\sum_{i \in I \cup A}$.

La separaci\'on en dos sumas es una consecuencia directa de tener dos conjuntos de alimentos con dominios distintos.
\end{cajazul}

%% =====================================================================
\section{Paso 5 --- Restricciones Nutricionales: M\'inimos que cumplir}
%% =====================================================================

Cada nutriente tiene un requerimiento m\'inimo diario. La dieta debe aportar \textbf{al menos} esa cantidad. La restricci\'on tiene la misma estructura que la funci\'on objetivo: dos sumatorias.

\subsection*{Ejemplo detallado: Prote\'ina ($\geq$ 65 g/d\'ia, $j = 1$)}

El aporte total de prote\'ina viene de \textbf{todos} los alimentos, tanto continuos como discretos:

\begin{align*}
    &\underbrace{3{,}3\,x_1 + 20{,}6\,x_3 + 6{,}4\,x_4 + 12{,}2\,x_5 + 0\,x_9}_{\text{aporte de alimentos continuos}} \\[4pt]
    + &\underbrace{13{,}5\,y_2 + 0{,}9\,y_6 + 0{,}3\,y_7 + 0{,}7\,y_8}_{\text{aporte de alimentos discretos}}
    \;\geq\; 65
\end{align*}

En notaci\'on compacta:
\[
    \sum_{i \in I} a_{i1} \cdot x_i + \sum_{i \in A} a_{i1} \cdot y_i \geq d_1
\]

\subsection*{Generalizaci\'on para los 8 nutrientes}

Cada nutriente $j \in J$ genera una restricci\'on con la misma estructura:

\begin{cajazul}
\[
    \sum_{i \in I} a_{ij} \cdot x_i + \sum_{i \in A} a_{ij} \cdot y_i \geq d_j, \quad \forall\, j \in J
\]

Esto dice: ``para \textbf{cada nutriente} $j$, la suma del aporte de todos los alimentos (continuos + discretos) debe ser al menos $d_j$''. El $\forall\, j \in J$ genera \textbf{8 restricciones} (una por nutriente).
\end{cajazul}

\begin{cajaverde}
\textbf{Compara con Panam\'a.} All\'a, $\forall\, j$ generaba restricciones \textit{por persona} (equipaje). Aqu\'i, $\forall\, j$ genera restricciones \textit{por nutriente}. El patr\'on es el mismo: cada vez que algo aplica ``para cada X'', usas $\forall$ sobre el conjunto correspondiente.
\end{cajaverde}

%% =====================================================================
\section{El Modelo Extendido (Expl\'icito)}
%% =====================================================================

El modelo extendido escribe cada restricci\'on con todos sus coeficientes, sin abreviar.

\begin{mdframed}[backgroundcolor=gray!5, linewidth=1.5pt]

\textbf{Funci\'on objetivo:}
\[
    \text{Min}\; z = 660\,x_1 + 2\,100\,x_3 + 950\,x_4 + 560\,x_5 + 800\,x_9
    + 1\,180\,y_2 + 440\,y_6 + 1\,440\,y_7 + 400\,y_8
\]

\textbf{Sujeto a:}

\medskip
\textit{Prote\'ina} ($\geq 65$ g):
\[
    3{,}3\,x_1 + 20{,}6\,x_3 + 6{,}4\,x_4 + 12{,}2\,x_5
    + 13{,}5\,y_2 + 0{,}9\,y_6 + 0{,}3\,y_7 + 0{,}7\,y_8 \geq 65
\]

\textit{L\'ipidos} ($\geq 66$ g):
\[
    3{,}2\,x_1 + 1{,}6\,x_3 + 0{,}8\,x_4 + 0{,}3\,x_5 + 82{,}9\,x_9
    + 10\,y_2 + 0{,}5\,y_6 + 0{,}3\,y_7 + 0{,}3\,y_8 \geq 66
\]

\textit{Hidratos de carbono} ($\geq 362$ g):
\[
    4{,}8\,x_1 + 57{,}3\,x_3 + 79{,}7\,x_4 + 74{,}6\,x_5
    + 4\,y_2 + 8{,}1\,y_6 + 14{,}5\,y_7 + 8{,}7\,y_8 \geq 362
\]

\textit{Vitamina A} ($\geq 700$ $\mu$g):
\[
    91\,x_1 + 85\,x_3 + 5{,}6\,x_9
    + 226\,y_2 + 1\,346\,y_6 + 4\,y_7 + 40\,y_8 \geq 700
\]

\textit{Vitamina C} ($\geq 75$ mg):
\[
    4{,}5\,x_1 + 26\,x_3
    + 6\,y_6 + 10\,y_7 + 50\,y_8 \geq 75
\]

\textit{Folato} ($\geq 400$ mg):
\[
    14\,x_1 + 88\,x_3 + 12\,x_4 + 4\,x_5 + 3\,x_9
    + 51{,}2\,y_2 + 10\,y_6 + 5\,y_7 + 37\,y_8 \geq 400
\]

\textit{Calcio} ($\geq 1\,000$ mg):
\[
    123\,x_1 + 260\,x_3 + 14\,x_4 + 24\,x_5
    + 100\,y_2 + 41\,y_6 + 6\,y_7 + 36\,y_8 \geq 1\,000
\]

\textit{Hierro} ($\geq 18$ mg):
\[
    0{,}1\,x_1 + 2\,x_3 + 0{,}8\,x_4 + 3{,}6\,x_5
    + 2{,}4\,y_2 + 0{,}7\,y_6 + 0{,}4\,y_7 + 0{,}3\,y_8 \geq 18
\]

\medskip
\textbf{Dominio:}
\begin{align*}
    x_i &\in \mathbb{R}^+ \cup \{0\}, \quad \forall\, i \in I = \{1, 3, 4, 5, 9\} \\
    y_i &\in \mathbb{Z}^+ \cup \{0\}, \quad \forall\, i \in A = \{2, 6, 7, 8\}
\end{align*}
\end{mdframed}

\begin{cajanota}
\textbf{\textquestiondown Cu\'antas variables y restricciones?}
\begin{itemize}
    \item \textbf{Variables continuas ($x_i$):} 5 (leche, porotos, arroz, fideos, mantequilla)
    \item \textbf{Variables enteras ($y_i$):} 4 (huevo, zanahoria, manzana, naranja)
    \item \textbf{Total variables:} 9
    \item \textbf{Restricciones nutricionales:} 8 (una por nutriente)
    \item \textbf{Tipo:} \textbf{PLEM} --- Programa Lineal Entero-Mixto de \textbf{minimizaci\'on}
\end{itemize}
\end{cajanota}

%% =====================================================================
\section{El Modelo Algebraico (Forma Compacta)}
%% =====================================================================

El modelo algebraico condensa todo lo anterior usando conjuntos, par\'ametros e \'indices.

\begin{mdframed}[backgroundcolor=gray!5, linewidth=1.5pt]
\textbf{Conjuntos:}
\begin{itemize}
    \item $I = \{1, 3, 4, 5, 9\}$ --- alimentos continuos (leche, porotos, arroz, fideos, mantequilla)
    \item $A = \{2, 6, 7, 8\}$ --- alimentos discretos (huevo, zanahoria, manzana, naranja)
    \item $J = \{1, 2, \ldots, 8\}$ --- nutrientes
\end{itemize}

\medskip
\textbf{Par\'ametros:}
\begin{itemize}
    \item $c_i$ : costo del alimento $i \in I \cup A$
    \item $a_{ij}$ : aporte nutricional del alimento $i \in I \cup A$ en el nutriente $j \in J$
    \item $d_j$ : requerimiento m\'inimo del nutriente $j \in J$
\end{itemize}

\medskip
\textbf{Variables de decisi\'on:}
\begin{itemize}
    \item $x_i$ : cantidad de porciones del alimento $i \in I$ (continua)
    \item $y_i$ : cantidad de porciones del alimento $i \in A$ (entera)
\end{itemize}

\bigskip
\textbf{Modelo:}
\begin{align}
    \text{Minimizar}\quad z &= \sum_{i \in I} c_i \cdot x_i + \sum_{i \in A} c_i \cdot y_i \tag{2.96} \\[10pt]
    \text{s.a.}\quad \sum_{i \in I} a_{ij} \cdot x_i + \sum_{i \in A} a_{ij} \cdot y_i &\geq d_j, \qquad \forall\, j \in J \tag{2.97} \\[10pt]
    x_i &\in \mathbb{R}^+ \cup \{0\}, \qquad \forall\, i \in I \tag{2.98} \\[10pt]
    y_i &\in \mathbb{Z}^+ \cup \{0\}, \qquad \forall\, i \in A \tag{2.99}
\end{align}
\end{mdframed}

\begin{cajazul}
\textbf{Interpretaci\'on de cada expresi\'on:}
\begin{itemize}
    \item \textbf{(2.96)} es la funci\'on objetivo: minimiza los costos asociados a la compra de alimentos. Las variables $x$ e $y$ surgen debido a las unidades de medida de los alimentos, siendo algunos continuos y otros discretos.

    \item \textbf{(2.97)} garantiza que la dieta aporte diariamente el m\'inimo de cada nutriente: prote\'ina, l\'ipidos, hidratos de carbono, vitamina A, vitamina C, folato, calcio y hierro.

    \item \textbf{(2.98)} y \textbf{(2.99)} representan la no-negatividad de las variables de decisi\'on: continuas para los alimentos medidos en tazas/cucharadas, y enteras para los alimentos que se compran por unidad.
\end{itemize}
\end{cajazul}

%% =====================================================================
\section{C\'odigo AMPL}
%% =====================================================================

El modelo algebraico se implementa directamente en AMPL (archivo \texttt{ProblemaDeLaDieta.mod}):

\begin{lstlisting}[style=ampl]
set I:={1, 3, 4, 5, 9};
set A:={2, 6, 7, 8};
set J:={1..8};
param c{I union A};
param d{J};
param a{I union A,J};
var x{I} >= 0;
var y{A} >=0 integer;
minimize z:sum{i in I} c[i]*x[i] + sum{i in A} c[i]*y[i];
s.t.
r1{j in J}: sum{i in I} a[i,j]*x[i] + sum{i in A} a[i,j]*y[i] >= d[j];
\end{lstlisting}

\begin{cajanota}
\textbf{Observa la correspondencia directa} entre el modelo algebraico y el c\'odigo AMPL:
\begin{itemize}
    \item Los conjuntos $I$, $A$, $J$ se declaran con \texttt{set}.
    \item Los par\'ametros $c_i$, $d_j$, $a_{ij}$ se declaran con \texttt{param}.
    \item \texttt{var x\{I\} >= 0} define las variables continuas.
    \item \texttt{var y\{A\} >= 0 integer} define las variables enteras.
    \item \texttt{sum\{i in I\}} corresponde exactamente a $\sum_{i \in I}$.
    \item \texttt{r1\{j in J\}} genera la restricci\'on $\forall\, j \in J$.
\end{itemize}
\end{cajanota}

%% =====================================================================
\section{Soluci\'on \'Optima (CPLEX)}
%% =====================================================================

Resolviendo con CPLEX en AMPL, la soluci\'on \'optima tiene un \textbf{costo m\'inimo diario de 9\,794,85 COP/d}. La dieta \'optima es:

\begin{table}[h]
\centering
\begin{tabular}{llr}
\toprule
\textbf{Variable} & \textbf{Alimento (porci\'on)} & \textbf{Valor \'optimo} \\
\midrule
$x_1$ & Leche entera ($\frac{1}{2}$ taza) & 4,071 \\
$x_3$ & Porotos (2 tazas) & 0,022 \\
$x_4$ & Arroz (2 tazas) & 0 \\
$x_5$ & Fideos (2 tazas) & 3,846 \\
$x_9$ & Mantequilla (4 cucharadas) & 0,582 \\
\midrule
$y_2$ & Huevo (2 unidades) & 0 \\
$y_6$ & Zanahoria (1 unidad) & 1 \\
$y_7$ & Manzana (1 unidad) & 0 \\
$y_8$ & Naranja (1 unidad) & 10 \\
\bottomrule
\end{tabular}
\caption{Soluci\'on \'optima: cantidad de alimentos m\'inimos para sobrevivir (Tabla 2.11 del libro).}
\end{table}

\begin{cajaverde}
\textbf{Interpretaci\'on de la soluci\'on:}
\begin{itemize}
    \item La dieta m\'as barata cuesta \textbf{9\,794,85 COP al d\'ia} (aproximadamente USD \$2,5).
    \item Las variables continuas toman valores fraccionarios: 4,071 medias-tazas de leche, 3,846 porciones de fideos, etc. Esto tiene sentido f\'isico (puedes medir fracciones de taza).
    \item Las variables enteras toman valores enteros: exactamente 1 zanahoria y 10 naranjas. No hay fracciones porque se declar\'o $y_i \in \mathbb{Z}^+$.
    \item \textbf{Arroz, huevo y manzana no se consumen} ($= 0$). Son demasiado caros para su aporte nutricional en comparaci\'on con las alternativas.
    \item \textbf{10 naranjas} parecen muchas, pero a 400 COP cada una son la fuente m\'as barata de Vitamina C y Folato.
\end{itemize}
\end{cajaverde}

%% =====================================================================
\section{Comparaci\'on Final: Dieta vs.\ Panam\'a}
%% =====================================================================

\begin{table}[h]
\centering
\small
\begin{tabular}{lll}
\toprule
\textbf{Aspecto} & \textbf{Panam\'a / Viajeros} & \textbf{Dieta} \\
\midrule
Objetivo              & Maximizar ganancia       & Minimizar costo \\
Sentido restricciones & $\leq$ (no exceder)      & $\geq$ (cumplir m\'inimos) \\
Conjuntos             & Personas $J$             & $I$ (continuos), $A$ (discretos), $J$ (nutrientes) \\
Variables             & Una sola ($x_{ij} \in \mathbb{Z}_{\geq 0}$) & Dos tipos ($x_i \in \mathbb{R}_{\geq 0}$, $y_i \in \mathbb{Z}_{\geq 0}$) \\
$\forall$ genera\ldots & restr.\ por persona     & restr.\ por nutriente \\
Tipo de problema      & PLE (entero puro)        & \textbf{PLEM} (entero-mixto) \\
\bottomrule
\end{tabular}
\caption{Resumen comparativo de ambos problemas.}
\end{table}

%% =====================================================================
\section{Resumen de Lecciones de Modelado}
%% =====================================================================

\begin{enumerate}[label=\textbf{\arabic*.}]
    \item \textbf{Analiza el dominio de cada variable antes de modelar.} No todos los alimentos se miden igual. Unos son divisibles (tazas) y otros no (unidades). Esto determina si la variable es continua o entera, y cambia el tipo de problema (PL $\to$ PLEM).

    \item \textbf{Separar conjuntos seg\'un el dominio.} Cuando hay variables de distinto tipo, crea conjuntos separados ($I$ para continuas, $A$ para discretas). Las sumatorias se separan en consecuencia.

    \item \textbf{El sentido de la optimizaci\'on determina el sentido de las restricciones.} Si minimizas costo, las restricciones t\'ipicamente son $\geq$ (cumplir m\'inimos). Si maximizas ganancia, t\'ipicamente son $\leq$ (no exceder recursos).

    \item \textbf{$\sum$ agrega dentro de una restricci\'on; $\forall$ replica restricciones.} En la dieta, $\sum_i$ suma el aporte de todos los alimentos; $\forall\, j$ genera una restricci\'on por nutriente.

    \item \textbf{El modelo algebraico se traduce directamente a AMPL.} Cada $\sum$ se convierte en \texttt{sum\{\}}, cada $\forall$ en un \'indice de restricci\'on \texttt{r\{j in J\}}, cada dominio en \texttt{>= 0} o \texttt{>= 0 integer}.
\end{enumerate}

\end{document}
